\chapter{Customizing Your Subversion
Experience}\label{svn.customization}

Version control can be a complex subject, as much art as science, that
offers myriad ways of getting stuff done. Throughout this book,
you\textquotesingle ve read of the various Subversion command-line
client subcommands and the options that modify their behavior. In this
chapter, we\textquotesingle ll look into still more ways to customize
the way Subversion works for you---setting up the Subversion runtime
configuration, using external helper applications,
Subversion\textquotesingle s interaction with the operating
system\textquotesingle s configured locale, and so on.

\section{Runtime Configuration Area}\label{svn.advanced.confarea}

Subversion provides many optional behaviors that the user can control.
Many of these options are of the kind that a user would wish to apply to
all Subversion operations. So, rather than forcing users to remember
command-line arguments for specifying these options and to use them for
every operation they perform, Subversion uses configuration files,
segregated into a Subversion configuration area.

The Subversion runtime configuration area is a two-tiered hierarchy of
option names and their values. Usually, this boils down to a special
directory that contains configuration files (the first tier), which are
just text files in standard INI format where ``sections'' provide the
second tier. You can easily edit these files using your favorite text
editor (such as Emacs or vi), and they contain directives read by the
client to determine which of several optional behaviors the user
prefers.

\subsection{Configuration Area
Layout}\label{svn.advanced.confarea.layout}

The first time the \texttt{svn} command-line client is executed, it
creates a per-user configuration area. On Unix-like systems, this area
appears as a directory named \texttt{.subversion} in the
user\textquotesingle s home directory. On Win32 systems, Subversion
creates a folder named \texttt{Subversion}, typically inside the
\texttt{Application\ Data} area of the user\textquotesingle s profile
directory (which, by the way, is usually a hidden directory). However,
on this platform, the exact location differs from system to system and
is dictated by the Windows Registry.\footnote{The \texttt{APPDATA}
  environment variable points to the \texttt{Application\ Data} area, so
  you can always refer to this folder as
  \texttt{\%APPDATA\%\textbackslash{}Subversion}.} We will refer to the
per-user configuration area using its Unix name, \texttt{.subversion}.

In addition to the per-user configuration area, Subversion also
recognizes the existence of a system-wide configuration area. This gives
system administrators the ability to establish defaults for all users on
a given machine. Note that the system-wide configuration area alone does
not dictate mandatory policy---the settings in the per-user
configuration area override those in the system-wide one, and
command-line arguments supplied to the \texttt{svn} program have the
final word on behavior. On Unix-like platforms, the system-wide
configuration area is expected to be the \texttt{/etc/subversion}
directory; on Windows machines, it looks for a \texttt{Subversion}
directory inside the common \texttt{Application\ Data} location (again,
as specified by the Windows Registry). Unlike the per-user case, the
\texttt{svn} program does not attempt to create the system-wide
configuration area.

The per-user configuration area currently contains three files---two
configuration files (\texttt{config} and \texttt{servers}), and a
\texttt{README.txt} file, which describes the INI format. At the time of
their creation, the files contain default values for each of the
supported Subversion options, mostly commented out and grouped with
textual descriptions about how the values for the key affect
Subversion\textquotesingle s behavior. To change a certain behavior, you
need only to load the appropriate configuration file into a text editor,
and to modify the desired option\textquotesingle s value. If at any time
you wish to have the default configuration settings restored, you can
simply remove (or rename) your configuration directory and then run some
innocuous \texttt{svn} command, such as \texttt{svn\ -\/-version}. A new
configuration directory with the default contents will be created.

Subversion also allows you to override individual configuration option
values at the command line via the \texttt{-\/-config-option} option,
which is especially useful if you need to make a (very) temporary change
in behavior. For more about this option\textquotesingle s proper usage,
see \hyperref[svn.ref.svn.sw]{svn Options}.

The per-user configuration area also contains a cache of authentication
data. The \texttt{auth} directory holds a set of subdirectories that
contain pieces of cached information used by
Subversion\textquotesingle s various supported authentication methods.
This directory is created in such a way that only the user herself has
permission to read its contents.

\subsection{Configuration and the Windows
Registry}\label{svn.advanced.confarea.windows-registry}

In addition to the usual INI-based configuration area, Subversion
clients running on Windows platforms may also use the Windows Registry
to hold the configuration data. The option names and their values are
the same as in the INI files. The ``file/section'' hierarchy is
preserved as well, though addressed in a slightly different fashion---in
this schema, files and sections are just levels in the Registry key
tree.

Subversion looks for system-wide configuration values under the
\texttt{HKEY\_LOCAL\_MACHINE\textbackslash{}Software\textbackslash{}Tigris.org\textbackslash{}Subversion}
key. For example, the \texttt{global-ignores} option, which is in the
\texttt{miscellany} section of the \texttt{config} file, would be found
at
\texttt{HKEY\_LOCAL\_MACHINE\textbackslash{}Software\textbackslash{}Tigris.org\textbackslash{}Subversion\textbackslash{}Config\textbackslash{}Miscellany\textbackslash{}global-ignores}.
Per-user configuration values should be stored under
\texttt{HKEY\_CURRENT\_USER\textbackslash{}Software\textbackslash{}Tigris.org\textbackslash{}Subversion}.

Registry-based configuration options are parsed \emph{before} their
file-based counterparts, so they are overridden by values found in the
configuration files. In other words, Subversion looks for configuration
information in the following locations on a Windows system;
lower-numbered locations take precedence over higher-numbered locations:

\begin{enumerate}
\def\labelenumi{\arabic{enumi}.}
\item
  Command-line options
\item
  The per-user INI files
\item
  The per-user Registry values
\item
  The system-wide INI files
\item
  The system-wide Registry values
\end{enumerate}

Also, the Windows Registry doesn\textquotesingle t really support the
notion of something being ``commented out.'' However, Subversion will
ignore any option key whose name begins with a hash (\texttt{\#})
character. This allows you to effectively comment out a Subversion
option without deleting the entire key from the Registry, obviously
simplifying the process of restoring that option.

The \texttt{svn} command-line client never attempts to write to the
Windows Registry and will not attempt to create a default configuration
area there. You can create the keys you need using the \texttt{REGEDIT}
program. Alternatively, you can create a \texttt{.reg} file (such as the
one in \hyperref[svn.advanced.confarea.windows-registry.ex-1]{Sample
registration entries (.reg) file}), and then double-click on that
file\textquotesingle s icon in the Explorer shell, which will cause the
data to be merged into your Registry.

\protect\phantomsection\label{svn.advanced.confarea.windows-registry.ex-1}
Sample registration entries (.reg) file

\begin{verbatim}
REGEDIT4

[HKEY_LOCAL_MACHINE\Software\Tigris.org\Subversion\Servers\groups]

[HKEY_LOCAL_MACHINE\Software\Tigris.org\Subversion\Servers\global]
"#http-auth-types"="basic;digest;negotiate"
"#http-compression"="yes"
"#http-library"=""
"#http-proxy-exceptions"=""
"#http-proxy-host"=""
"#http-proxy-password"=""
"#http-proxy-port"=""
"#http-proxy-username"=""
"#http-timeout"="0"
"#neon-debug-mask"=""
"#ssl-authority-files"=""
"#ssl-client-cert-file"=""
"#ssl-client-cert-password"=""
"#ssl-pkcs11-provider"=""
"#ssl-trust-default-ca"=""
"#store-auth-creds"="yes"
"#store-passwords"="yes"
"#store-plaintext-passwords"="ask"
"#store-ssl-client-cert-pp"="yes"
"#store-ssl-client-cert-pp-plaintext"="ask"
"#username"=""

[HKEY_CURRENT_USER\Software\Tigris.org\Subversion\Config\auth]
"#password-stores"="windows-cryptoapi"

[HKEY_CURRENT_USER\Software\Tigris.org\Subversion\Config\helpers]
"#diff-cmd"=""
"#diff-extensions"="-u"
"#diff3-cmd"=""
"#diff3-has-program-arg"=""
"#editor-cmd"="notepad"
"#merge-tool-cmd"=""

[HKEY_CURRENT_USER\Software\Tigris.org\Subversion\Config\tunnels]

[HKEY_CURRENT_USER\Software\Tigris.org\Subversion\Config\miscellany]
"#enable-auto-props"="no"
"#global-ignores"="*.o *.lo *.la *.al .libs *.so *.so.[0-9]* *.a *.pyc *.pyo *.rej *~ #*# .#* .*.swp .DS_Store"
"#interactive-conflicts"="yes"
"#log-encoding"=""
"#mime-types-file"=""
"#no-unlock"="no"
"#preserved-conflict-file-exts"="doc ppt xls od?"
"#use-commit-times"="no"

[HKEY_CURRENT_USER\Software\Tigris.org\Subversion\Config\auto-props]
\end{verbatim}

\hyperref[svn.advanced.confarea.windows-registry.ex-1]{Sample
registration entries (.reg) file} shows the contents of a \texttt{.reg}
file, which contains some of the most commonly used configuration
options and their default values. Note the presence of both system-wide
(for network proxy-related options) and per-user settings (editor
programs and password storage, among others). Also note that all the
options are effectively commented out. You need only to remove the hash
(\texttt{\#}) character from the beginning of the option names and set
the values as you desire.

\subsection{Runtime Configuration
Options}\label{svn.advanced.confarea.opts}

In this section, we will discuss the specific runtime configuration
options that Subversion currently supports.

\subsubsection{General
configuration}\label{svn.advanced.confarea.opts.config}

The \texttt{config} file contains the rest of the currently available
Subversion runtime options---those not related to networking. There are
only a few options in use as of this writing, but they are again grouped
into sections in expectation of future additions.

The \texttt{{[}auth{]}} section contains settings related to
Subversion\textquotesingle s authentication and authorization against
the repository. It contains the following:

\begin{description}
\item[\texttt{password-stores}]
This comma-delimited list specifies which (if any) system-provided
password stores Subversion should attempt to use when saving and
retrieving cached authentication credentials, and in what order
Subversion should prefer them. The default value is
\texttt{gnome-keyring,\ kwallet,\ keychain,\ gpg-agent,\ windows-crypto-api},
representing the GNOME Keyring, KDE Wallet, Mac OS X Keychain, GnuPG
Agent, and Microsoft Windows cryptography API, respectively. Listed
stores which are not available on the system are ignored.
\item[\texttt{store-passwords}]
This option has been deprecated from the \texttt{config} file. It now
lives as a per-server configuration item in the \texttt{servers}
configuration area. See
\hyperref[svn.advanced.confarea.opts.servers]{Per-server configuration}
for details.
\item[\texttt{store-auth-creds}]
This option has been deprecated from the \texttt{config} file. It now
lives as a per-server configuration item in the \texttt{servers}
configuration area. See
\hyperref[svn.advanced.confarea.opts.servers]{Per-server configuration}
for details.
\end{description}

The \texttt{{[}helpers{]}} section controls which external applications
Subversion uses to accomplish its tasks. Valid options in this section
are:

\begin{description}
\item[\texttt{diff-cmd}]
This specifies the absolute path of a differencing program, used when
Subversion generates ``diff'' output (such as when using the
\texttt{svn\ diff} command). By default, Subversion uses an internal
differencing library---setting this option will cause it to perform this
task using an external program. See
\hyperref[svn.advanced.externaldifftools]{Using External Differencing
and Merge Tools} for more details on using such programs.
\item[\texttt{diff-extensions}]
Like the \texttt{-\/-extensions} (\texttt{-x}) command-line option, this
specifies additional options passed to the file content differencing
engine. The set of meaningful extension options differs depending on
whether the client is using Subversion\textquotesingle s internal
differencing engine or an external mechanism. See the output of
\texttt{svn\ help\ diff} for details. The default value for this option
is \texttt{-u}.
\item[\texttt{diff3-cmd}]
This specifies the absolute path of a three-way differencing program.
Subversion uses this program to merge changes made by the user with
those received from the repository. By default, Subversion uses an
internal differencing library---setting this option will cause it to
perform this task using an external program. See
\hyperref[svn.advanced.externaldifftools]{Using External Differencing
and Merge Tools} for more details on using such programs.
\item[\texttt{diff3-has-program-arg}]
This flag should be set to \texttt{true} if the program specified by the
\texttt{diff3-cmd} option accepts a \texttt{-\/-diff-program}
command-line parameter.
\item[\texttt{editor-cmd}]
This specifies the program Subversion will use to query the user for
certain types of textual metadata or when interactively resolving
conflicts. See \hyperref[svn.advanced.externaleditors]{Using External
Editors} for more details on using external text editors with
Subversion.
\item[\texttt{merge-tool-cmd}]
This specifies the program that Subversion will use to perform three-way
merge operations on your versioned files. See
\hyperref[svn.advanced.externaldifftools]{Using External Differencing
and Merge Tools} for more details on using such programs.
\end{description}

The \texttt{{[}tunnels{]}} section allows you to define new tunnel
schemes for use with \texttt{svnserve} and \texttt{svn://} client
connections. For more details, see
\hyperref[svn.serverconfig.svnserve.sshauth]{Tunneling over SSH}.

The \texttt{{[}miscellany{]}} section is where everything that
doesn\textquotesingle t belong elsewhere winds up.\footnote{Anyone for
  potluck dinner?} In this section, you can find:

\begin{description}
\item[\texttt{enable-auto-props}]
This instructs Subversion to automatically set properties on newly added
or imported files. The default value is \texttt{no}, so set this to
\texttt{yes} to enable this feature. The \texttt{{[}auto-props{]}}
section of this file specifies which properties are to be set on which
files.
\item[\texttt{global-ignores}]
When running the \texttt{svn\ status} command, Subversion lists
unversioned files and directories along with the versioned ones,
annotating them with a \texttt{?} character (see
\hyperref[svn.tour.cycle.examine.status]{See an overview of your
changes}). Sometimes it can be annoying to see uninteresting,
unversioned items---for example, object files that result from a
program\textquotesingle s compilation---in this display. The
\texttt{global-ignores} option is a list of whitespace-delimited globs
that describe the names of files and directories that Subversion should
not display unless they are versioned. The default value is
\texttt{*.o\ *.lo\ *.la\ *.al\ .libs\ *.so\ *.so.{[}0-9{]}*\ *.a\ *.pyc\ *.pyo\ *.rej\ *\textasciitilde{}\ \#*\#\ .\#*\ .*.swp\ .DS\_Store}.

As well as \texttt{svn\ status}, the \texttt{svn\ add} and
\texttt{svn\ import} commands also ignore files that match the list when
they are scanning a directory. You can override this behavior for a
single instance of any of these commands by explicitly specifying the
filename, or by using the \texttt{-\/-no-ignore} command-line flag.

For information on finer-grained control of ignored items, see
\hyperref[svn.advanced.props.special.ignore]{Ignoring Unversioned
Items}.
\item[\texttt{interactive-conflicts}]
This is a Boolean option that specifies whether Subversion should try to
resolve conflicts interactively. If its value is \texttt{yes} (which is
the default value), Subversion will prompt the user for how to handle
conflicts in the manner demonstrated in
\hyperref[svn.tour.cycle.resolve]{Resolve Any Conflicts}. Otherwise, it
will simply flag the conflict and continue its operation, postponing
resolution to a later time.
\item[\texttt{log-encoding}]
This variable sets the default character set encoding for commit log
messages. It\textquotesingle s a permanent form of the
\texttt{-\/-encoding} option (see \hyperref[svn.ref.svn.sw]{svn
Options}). The Subversion repository stores log messages in UTF-8 and
assumes that your log message is written using your operating
system\textquotesingle s native locale. You should specify a different
encoding if your commit messages are written in any other encoding.
\item[\texttt{mime-types-file}]
This option, new to Subversion 1.5, specifies the path of a MIME types
mapping file, such as the \texttt{mime.types} file provided by the
Apache HTTP Server. Subversion uses this file to assign MIME types to
newly added or imported files. See
\hyperref[svn.advanced.props.auto]{Automatic Property Setting} and
\hyperref[svn.advanced.props.special.mime-type]{File Content Type} for
more about Subversion\textquotesingle s detection and use of file
content types.
\item[\texttt{no-unlock}]
This Boolean option corresponds to
\texttt{svn\ commit}\textquotesingle s \texttt{-\/-no-unlock} option,
which tells Subversion not to release locks on files
you\textquotesingle ve just committed. If this runtime option is set to
\texttt{yes}, Subversion will never release locks automatically, leaving
you to run \texttt{svn\ unlock} explicitly. It defaults to \texttt{no}.
\item[\texttt{preserved-conflict-file-exts}]
The value of this option is a space-delimited list of file extensions
that Subversion should preserve when generating conflict filenames. By
default, the list is empty. This option is new to Subversion 1.5.

When Subversion detects conflicting file content changes, it defers
resolution of those conflicts to the user. To assist in the resolution,
Subversion keeps pristine copies of the various competing versions of
the file in the working copy. By default, those conflict files have
names constructed by appending to the original filename a custom
extension such as \texttt{.mine} or \texttt{.REV} (where
\textless REV\textgreater{} is a revision number). A mild annoyance with
this naming scheme is that on operating systems where a
file\textquotesingle s extension determines the default application used
to open and edit that file, appending a custom extension prevents the
file from being easily opened by its native application. For example, if
the file \texttt{ReleaseNotes.pdf} was conflicted, the conflict files
might be named \texttt{ReleaseNotes.pdf.mine} or
\texttt{ReleaseNotes.pdf.r4231}. While your system might be configured
to use Adobe\textquotesingle s Acrobat Reader to open files whose
extensions are \texttt{.pdf}, there probably isn\textquotesingle t an
application configured on your system to open all files whose extensions
are \texttt{.r4231}.

You can fix this annoyance by using this configuration option, though.
For files with one of the specified extensions, Subversion will append
to the conflict file names the custom extension just as before, but then
also reappend the file\textquotesingle s original extension. Using the
previous example, and assuming that \texttt{pdf} is one of the
extensions configured in this list thereof, the conflict files generated
for \texttt{ReleaseNotes.pdf} would instead be named
\texttt{ReleaseNotes.pdf.mine.pdf} and
\texttt{ReleaseNotes.pdf.r4231.pdf}. Because each file ends in
\texttt{.pdf}, the correct default application will be used to view
them.
\item[\texttt{use-commit-times}]
Normally your working copy files have timestamps that reflect the last
time they were touched by any process, whether your own editor or some
\texttt{svn} subcommand. This is generally convenient for people
developing software, because build systems often look at timestamps as a
way of deciding which files need to be recompiled.

In other situations, however, it\textquotesingle s sometimes nice for
the working copy files to have timestamps that reflect the last time
they were changed in the repository. The \texttt{svn\ export} command
always places these ``last-commit timestamps'' on trees that it
produces. By setting this config variable to \texttt{yes}, the
\texttt{svn\ checkout}, \texttt{svn\ update}, \texttt{svn\ switch}, and
\texttt{svn\ revert} commands will also set last-commit timestamps on
files that they touch.
\end{description}

The \texttt{{[}auto-props{]}} section controls the Subversion
client\textquotesingle s ability to automatically set properties on
files when they are added or imported. It contains any number of
key-value pairs in the format
\texttt{PATTERN\ =\ PROPNAME=VALUE{[};PROPNAME=VALUE\ ...{]}}, where
\textless PATTERN\textgreater{} is a file pattern that matches one or
more filenames and the rest of the line is a semicolon-delimited set of
property assignments. (If you need to use a semicolon in your
property\textquotesingle s name or value, you can escape it by doubling
it.)

\begin{verbatim}
$ cat ~/.subversion/config
…
[auto-props]
*.c = svn:eol-style=native
*.html = svn:eol-style=native;svn:mime-type=text/html;; charset=UTF8
*.sh = svn:eol-style=native;svn:executable
…
$ cd projects/myproject
$ svn status
?       www/index.html
$ svn add www/index.html
A         www/index.html
$ svn diff www/index.html
…

Property changes on: www/index.html
___________________________________________________________________
Added: svn:mime-type
## -0,0 +1 ##
+text/html; charset=UTF8
Added: svn:eol-style
## -0,0 +1 ##
+native
$
\end{verbatim}

Multiple matches on a file will result in multiple propsets for that
file; however, there is no guarantee that auto-props will be applied in
the order in which they are listed in the config file, so you
can\textquotesingle t have one rule ``override'' another. You can find
several examples of auto-props usage in the \texttt{config} file.
Lastly, don\textquotesingle t forget to set \texttt{enable-auto-props}
to \texttt{yes} in the \texttt{{[}miscellany{]}} section if you want to
enable auto-props.

New to Subversion 1.8, the \texttt{{[}working-copy{]}} section is used
for configuring working copies. Valid options in this section are:

\begin{description}
\item[\texttt{exclusive-locking-clients}]
Enables exclusive SQLite locking of working copies for the client, hence
improving performance for working copies located on network disks. By
setting this config variable to \texttt{svn}, you instruct Subversion
command-line client to use exclusive locking. This reduces the locking
overhead but does mean that only one Subversion client will be able to
access the working copy at a time. A second client attempting to access
a locked working copy will block for 10 seconds and then get an error.
In most cases shared locking is preferred but if the working copy is on
a network disk rather than a local disk the locking overhead is more
significant. When dealing with large working copies on network disks
exclusive locking may give a significant performance gain, two or three
times faster in some cases. This option is new to Subversion 1.8.
\item[\texttt{exclusive-locking}]
Setting this config variable to \texttt{true} enables exclusive SQLite
locking of working copies for all Subversion 1.8 clients. Enabling this
may cause some clients to fail to work properly. The default value for
this option is \texttt{false}. This option is new to Subversion 1.8.
\end{description}

\subsubsection{Per-server
configuration}\label{svn.advanced.confarea.opts.servers}

The \texttt{servers} file contains Subversion configuration options
related to the network layers. There are two special sections in this
file---\texttt{{[}groups{]}} and \texttt{{[}global{]}}. The
\texttt{{[}groups{]}} section is essentially a cross-reference table.
The keys in this section are the names of other sections in the file;
their values are globs---textual tokens that possibly contain wildcard
characters---that are compared against the hostnames of the machine to
which Subversion requests are sent.

\begin{verbatim}
[groups]
beanie-babies = *.red-bean.com
collabnet = svn.collab.net

[beanie-babies]
…

[collabnet]
…
\end{verbatim}

When Subversion is used over a network, it attempts to match the name of
the server it is trying to reach with a group name under the
\texttt{{[}groups{]}} section. If a match is made, Subversion then looks
for a section in the \texttt{servers} file whose name is the matched
group\textquotesingle s name. From that section, it reads the actual
network configuration settings.

The \texttt{{[}global{]}} section contains the settings that are meant
for all of the servers not matched by one of the globs under the
\texttt{{[}groups{]}} section. The options available in this section are
exactly the same as those that are valid for the other server sections
in the file (except, of course, the special \texttt{{[}groups{]}}
section), and are as follows:

\begin{description}
\item[\texttt{http-auth-types}]
This is a semicolon-delimited list of HTTP authentication types which
the client will deem acceptable. Valid types are \texttt{basic},
\texttt{digest}, and \texttt{negotiate}, with the default behavior being
acceptance of any these authentication types. A client which insists on
not transmitting authentication credentials in cleartext might, for
example, be configured such that the value of this option is
\texttt{digest;negotiate}---omitting \texttt{basic} from the list. (Note
that this setting is only honored by Subversion\textquotesingle s
Neon-based HTTP provider module.)
\item[\texttt{http-compression}]
This specifies whether Subversion should attempt to compress network
requests made to DAV-ready servers. The default value is \texttt{yes}
(though compression will occur only if that capability is compiled into
the network layer). Set this to \texttt{no} to disable compression, such
as when debugging network transmissions.
\item[\texttt{http-library}]
The \texttt{http-library} runtime configuration option allows users to
specify (generally, or in a per-server-group fashion) which of the
available WebDAV access modules they\textquotesingle d prefer to use.
Prior to version 1.8, Subversion offered a pair of such modules: its
original implementiation \texttt{libsvn\_ra\_neon} (selected by using
the value \texttt{neon} for this option) and the newer
\texttt{libsvn\_ra\_serf} (selected using the value \texttt{serf}). As
of Subversion 1.8, only \texttt{libsvn\_ra\_serf} is supported. This
configuration option remains, though, because the runtime configuration
area is version-agnostic. Users with multiple versions of Subversion
installed may still wish to enable the use of \texttt{libsvn\_ra\_neon}
for sites which they access with an older version of Subversion.
\item[\texttt{http-proxy-exceptions}]
This specifies a comma-separated list of patterns for repository
hostnames that should be accessed directly, without using the proxy
machine. The pattern syntax is the same as is used in the Unix shell for
filenames. A repository hostname matching any of these patterns will not
be proxied.
\item[\texttt{http-proxy-host}]
This specifies the hostname of the proxy computer through which your
HTTP-based Subversion requests must pass. It defaults to an empty value,
which means that Subversion will not attempt to route HTTP requests
through a proxy computer, and will instead attempt to contact the
destination machine directly.
\item[\texttt{http-proxy-password}]
This specifies the password to supply to the proxy machine. It defaults
to an empty value.
\item[\texttt{http-proxy-port}]
This specifies the port number on the proxy host to use. It defaults to
an empty value.
\item[\texttt{http-proxy-username}]
This specifies the username to supply to the proxy machine. It defaults
to an empty value.
\item[\texttt{http-timeout}]
This specifies the amount of time, in seconds, to wait for a server
response. If you experience problems with a slow network connection
causing Subversion operations to time out, you should increase the value
of this option. In Subversion 1.8 (or older versions employing the
Serf-based HTTP provider), use the value \texttt{0} to disable the
timeout altogether.
\item[\texttt{neon-debug-mask}]
This is an integer mask that the Neon HTTP library uses for choosing
what type of debugging output to yield. The default value is \texttt{0},
which will silence all debugging output. Prior to version 1.8, most
Subversion clients used Neon (via the \texttt{libsvn\_ra\_neon}
repository access module) for WebDAV/HTTP communications between the
Subversion client and server. Support for \texttt{libsvn\_ra\_neon} was
dropped in Subversion 1.8, though, making this option obsolete for newer
Subversion installations.
\item[\texttt{ssl-authority-files}]
This is a semicolon-delimited list of paths to files containing
certificates of the certificate authorities (or CAs) that are accepted
by the Subversion client when accessing the repository over HTTPS.
\item[\texttt{ssl-client-cert-file}]
If a host (or set of hosts) requires an SSL client certificate,
you\textquotesingle ll normally be prompted for a path to your
certificate. By setting this variable to that same path, Subversion will
be able to find your client certificate automatically without prompting
you. There\textquotesingle s no standard place to store your certificate
on disk; Subversion will grab it from any path you specify.
\item[\texttt{ssl-client-cert-password}]
If your SSL client certificate file is encrypted by a passphrase,
Subversion will prompt you for the passphrase whenever the certificate
is used. If you find this annoying (and don\textquotesingle t mind
storing the password in the \texttt{servers} file), you can set this
variable to the certificate\textquotesingle s passphrase. You
won\textquotesingle t be prompted anymore.
\item[\texttt{ssl-pkcs11-provider}]
The value of this option is the name of the PKCS\#11 provider from which
an SSL client certificate will be drawn (if the server asks for one).
This setting is only honored by Subversion\textquotesingle s Neon-based
HTTP provider module, which was removed in Subversion 1.8.
\item[\texttt{ssl-trust-default-ca}]
Set this variable to \texttt{yes} if you want Subversion to
automatically trust the set of default CAs that ship with OpenSSL.
\item[\texttt{store-auth-creds}]
This setting is the same as \texttt{store-passwords}, except that it
enables or disables on-disk caching of \emph{all} authentication
information: usernames, passwords, server certificates, and any other
types of cacheable credentials.
\item[\texttt{store-passwords}]
This instructs Subversion to cache, or not to cache, passwords that are
supplied by the user in response to server authentication challenges.
The default value is \texttt{yes}. Set this to \texttt{no} to disable
this on-disk password caching. You can override this option for a single
instance of the \texttt{svn} command using the
\texttt{-\/-no-auth-cache} command-line parameter (for those subcommands
that support it). For more information regarding that, see
\hyperref[svn.serverconfig.netmodel.credcache]{Caching credentials}.
Note that regardless of how this option is configured, Subversion will
not store passwords in plaintext unless the
\texttt{store-plaintext-passwords} option is also set to \texttt{yes}.
\item[\texttt{store-plaintext-passwords}]
This variable is only important on UNIX-like systems. It controls what
the Subversion client does in case the password for the current
authentication realm can only be cached on disk in unencrypted form, in
the \texttt{\textasciitilde{}/.subversion/auth/} caching area. You can
set it to \texttt{yes} or \texttt{no} to enable or disable caching of
passwords in unencrypted form, respectively. The default setting is
\texttt{ask}, which causes the Subversion client to ask you each time a
\emph{new} password is about to be added to the
\texttt{\textasciitilde{}/.subversion/auth/} caching area.
\item[\texttt{store-ssl-client-cert-pp}]
This option controls whether Subversion will cache SSL client
certificate passphrases provided by the user. Its value defaults to
\texttt{yes}. Set this to \texttt{no} to disable this passphrase
caching.
\item[\texttt{store-ssl-client-cert-pp-plaintext}]
This option controls whether Subversion, when attempting to cache an SSL
client certificate passphrase, will be allowed to do so using its
on-disk plaintext storage mechanism. The default value of this option is
\texttt{ask}, which causes the Subversion client to ask you each time a
\emph{new} client certificate passphrase is about to be added to the
\texttt{\textasciitilde{}/.subversion/auth/} caching area. Set this
option\textquotesingle s value to \texttt{yes} or \texttt{no} to
indicate your preference and avoid related prompts.
\end{description}

\section{Localization}\label{svn.advanced.l10n}

Localization is the act of making programs behave in a region-specific
way. When a program formats numbers or dates in a way specific to your
part of the world or prints messages (or accepts input) in your native
language, the program is said to be localized. This section describes
steps Subversion has made toward localization.

\subsection{Understanding
Locales}\label{svn.advanced.l10n.understanding}

Most modern operating systems have a notion of the ``current
locale''---that is, the region or country whose localization conventions
are honored. These conventions---typically chosen by some runtime
configuration mechanism on the computer---affect the way in which
programs present data to the user, as well as the way in which they
accept user input.

On most Unix-like systems, you can check the values of the
locale-related runtime configuration options by running the
\texttt{locale} command:

\begin{verbatim}
$ locale
LANG=
LC_COLLATE="C"
LC_CTYPE="C"
LC_MESSAGES="C"
LC_MONETARY="C"
LC_NUMERIC="C"
LC_TIME="C"
LC_ALL="C"
$
\end{verbatim}

The output is a list of locale-related environment variables and their
current values. In this example, the variables are all set to the
default \texttt{C} locale, but users can set these variables to specific
country/language code combinations. For example, if one were to set the
\texttt{LC\_TIME} variable to \texttt{fr\_CA}, programs would know to
present time and date information formatted according to a
French-speaking Canadian\textquotesingle s expectations. And if one were
to set the \texttt{LC\_MESSAGES} variable to \texttt{zh\_TW}, programs
would know to present human-readable messages in Traditional Chinese.
Setting the \texttt{LC\_ALL} variable has the effect of changing every
locale variable to the same value. The value of \texttt{LANG} is used as
a default value for any locale variable that is unset. To see the list
of available locales on a Unix system, run the command
\texttt{locale\ -a}.

On Windows, locale configuration is done via the ``Regional and Language
Options'' control panel item. There you can view and select the values
of individual settings from the available locales, and even customize
(at a sickening level of detail) several of the display formatting
conventions.

\subsection{Subversion\textquotesingle s Use of
Locales}\label{svn.advanced.l10n.svnuse}

The Subversion client, \texttt{svn}, honors the current locale
configuration in two ways. First, it notices the value of the
\texttt{LC\_MESSAGES} variable and attempts to print all messages in the
specified language. For example:

\begin{verbatim}
$ export LC_MESSAGES=de_DE
$ svn help cat
cat: Gibt den Inhalt der angegebenen Dateien oder URLs aus.
Aufruf: cat ZIEL[@REV]...
…
\end{verbatim}

This behavior works identically on both Unix and Windows systems. Note,
though, that while your operating system might have support for a
certain locale, the Subversion client still may not be able to speak the
particular language. In order to produce localized messages, human
volunteers must provide translations for each language. The translations
are written using the GNU gettext package, which results in translation
modules that end with the \texttt{.mo} filename extension. For example,
the German translation file is named \texttt{de.mo}. These translation
files are installed somewhere on your system. On Unix, they typically
live in \texttt{/usr/share/locale/}, while on Windows
they\textquotesingle re often found in the
\texttt{share\textbackslash{}locale\textbackslash{}} folder in
Subversion\textquotesingle s installation area. Once installed, a module
is named after the program for which it provides translations. For
example, the \texttt{de.mo} file may ultimately end up installed as
\texttt{/usr/share/locale/de/LC\_MESSAGES/subversion.mo}. By browsing
the installed \texttt{.mo} files, you can see which languages the
Subversion client is able to speak.

The second way in which the locale is honored involves how \texttt{svn}
interprets your input. The repository stores all paths, filenames, and
log messages in Unicode, encoded as UTF-8. In that sense, the repository
is internationalized---that is, the repository is ready to accept input
in any human language. This means, however, that the Subversion client
is responsible for sending only UTF-8 filenames and log messages into
the repository. To do this, it must convert the data from the native
locale into UTF-8.

For example, suppose you create a file named \texttt{caffè.txt}, and
then when committing the file, you write the log message as ``Adesso il
caffè è più forte.'' Both the filename and the log message contain
non-ASCII characters, but because your locale is set to \texttt{it\_IT},
the Subversion client knows to interpret them as Italian. It uses an
Italian character set to convert the data to UTF-8 before sending it off
to the repository.

Note that while the repository demands UTF-8 filenames and log messages,
it \emph{does not} pay attention to file contents. Subversion treats
file contents as opaque strings of bytes, and neither client nor server
makes an attempt to understand the character set or encoding of the
contents.

Character Set Conversion Errors

While using Subversion, you might get hit with an error related to
character set conversions:

\begin{verbatim}
svn: E000022: Can't convert string from native encoding to 'UTF-8':
…
svn: E000022: Can't convert string from 'UTF-8' to native encoding:
…
\end{verbatim}

Errors such as this typically occur when the Subversion client has
received a UTF-8 string from the repository, but not all of the
characters in that string can be represented using the encoding of the
current locale. For example, if your locale is \texttt{en\_US} but a
collaborator has committed a Japanese filename, you\textquotesingle re
likely to see this error when you receive the file during an
\texttt{svn\ update}.

The solution is either to set your locale to something that \emph{can}
represent the incoming UTF-8 data, or to change the filename or log
message in the repository. (And don\textquotesingle t forget to slap
your collaborator\textquotesingle s hand---projects should decide on
common languages ahead of time so that all participants are using the
same locale.)

\section{Using External Editors}\label{svn.advanced.externaleditors}

The most obvious way to get data into Subversion is through the addition
of files to version control, committing changes to those files, and so
on. But other pieces of information besides merely versioned file data
live in your Subversion repository. Some of these bits of
information---commit log messages, lock comments, and some property
values---tend to be textual in nature and are provided explicitly by
users. Most of this information can be provided to the Subversion
command-line client using the \texttt{-\/-message} (\texttt{-m}) and
\texttt{-\/-file} (\texttt{-F}) options with the appropriate
subcommands.

Each of these options has its pros and cons. For example, when
performing a commit, \texttt{-\/-file} (\texttt{-F}) works well if
you\textquotesingle ve already prepared a text file that holds your
commit log message. If you didn\textquotesingle t, though, you can use
\texttt{-\/-message} (\texttt{-m}) to provide a log message on the
command line. Unfortunately, it can be tricky to compose anything more
than a simple one-line message on the command line. Users want more
flexibility---multiline, free-form log message editing on demand.

Subversion supports this by allowing you to specify an external text
editor that it will launch as necessary to give you a more powerful
input mechanism for this textual metadata. There are several ways to
tell Subversion which editor you\textquotesingle d like use. Subversion
checks the following things, in the order specified, when it wants to
launch such an editor:

\begin{enumerate}
\def\labelenumi{\arabic{enumi}.}
\item
  \texttt{-\/-editor-cmd} command-line option
\item
  \texttt{SVN\_EDITOR} environment variable
\item
  \texttt{editor-cmd} runtime configuration option
\item
  \texttt{VISUAL} environment variable
\item
  \texttt{EDITOR} environment variable
\item
  Possibly, a fallback value built into the Subversion libraries (not
  present in the official builds)
\end{enumerate}

The value of any of these options or variables is the beginning of a
command line to be executed by the shell. Subversion appends to that
command line a space and the pathname of a temporary file to be edited.
So, to be used with Subversion, the configured or specified editor needs
to support an invocation in which its last command-line parameter is a
file to be edited, and it should be able to save the file in place and
return a zero exit code to indicate success.

As noted, external editors can be used to provide commit log messages to
any of the committing subcommands (such as \texttt{svn\ commit} or
\texttt{import}, \texttt{svn\ mkdir} or \texttt{delete} when provided a
URL target, etc.), and Subversion will try to launch the editor
automatically if you don\textquotesingle t specify either of the
\texttt{-\/-message} (\texttt{-m}) or \texttt{-\/-file} (\texttt{-F})
options. The \texttt{svn\ propedit} command is built almost entirely
around the use of an external editor. And beginning in version 1.5,
Subversion will also use the configured external text editor when the
user asks it to launch an editor during interactive conflict resolution.
Oddly, there doesn\textquotesingle t appear to be a way to use external
editors to interactively provide lock comments.

\section{Using External Differencing and Merge
Tools}\label{svn.advanced.externaldifftools}

The interface between Subversion and external two- and three-way
differencing tools harkens back to a time when
Subversion\textquotesingle s only contextual differencing capabilities
were built around invocations of the GNU diffutils toolchain,
specifically the \texttt{diff} and \texttt{diff3} utilities. To get the
kind of behavior Subversion needed, it called these utilities with more
than a handful of options and parameters, most of which were quite
specific to the utilities. Some time later, Subversion grew its own
internal differencing library, and as a failover mechanism, the
\texttt{-\/-diff-cmd} and \texttt{-\/-diff3-cmd} options were added to
the Subversion command-line client so that users could more easily
indicate that they preferred to use the GNU diff and diff3 utilities
instead of the newfangled internal diff library. If those options were
used, Subversion would simply ignore the internal diff library, and fall
back to running those external programs, lengthy argument lists and all.
And that\textquotesingle s where things remain today.

It didn\textquotesingle t take long for folks to realize that having
such easy configuration mechanisms for specifying that Subversion should
use the external GNU diff and diff3 utilities located at a particular
place on the system could be applied toward the use of other
differencing tools, too. After all, Subversion didn\textquotesingle t
actually verify that the things it was being told to run were members of
the GNU diffutils toolchain. But the only configurable aspect of using
those external tools is their location on the system---not the option
set, parameter order, and so on. Subversion continues to throw all those
GNU utility options at your external diff tool regardless of whether
that program can understand those options. And that\textquotesingle s
where things get unintuitive for most users.

The decision on when to fire off a contextual two- or three-way diff as
part of a larger Subversion operation is made internally by Subversion
and is affected by, among other things, whether the files being operated
on are human-readable as determined by their \texttt{svn:mime-type}
property. This means, for example, that even if you had the niftiest
Microsoft Word-aware differencing or merging tool in the universe, it
would typically not be invoked by Subversion if your versioned Word
documents had a configured MIME type that denoted that they were not
human-readable (such as \texttt{application/msword}). Fortunately, you
can pass the \texttt{-\/-force} option to \texttt{svn\ diff} to
short-circuit this MIME-related sanity check and force the difference to
be calculated. For more about MIME type settings, see
\hyperref[svn.advanced.props.special.mime-type]{File Content Type}

Much later, Subversion 1.5 introduced interactive resolution of
conflicts (described in \hyperref[svn.tour.cycle.resolve]{Resolve Any
Conflicts}). One of the options that this feature provides to users is
the ability to interactively launch a third-party merge tool. If this
action is taken, Subversion will check to see if the user has specified
such a tool for use in this way. Subversion will first check the
\texttt{SVN\_MERGE} environment variable for the name of an external
merge tool. If that variable is not set, it will look for the same
information in the value of the \texttt{merge-tool-cmd} runtime
configuration option. Upon finding a configured external merge tool, it
will invoke that tool.

While the general purposes of the three-way differencing and merge tools
are roughly the same (finding a way to make separate-but-overlapping
file changes live in harmony), Subversion exercises each of these
options at different times and for different reasons. The internal
three-way differencing engine and its optional external replacement are
used when interaction with the user is \emph{not} expected. In fact,
significant delay introduced by such a tool can actually result in the
failure of some time-sensitive Subversion operations.
It\textquotesingle s the external merge tool that is intended to be
invoked interactively.

Now, while the interface between Subversion and an external merge tool
is significantly less convoluted than that between Subversion and the
diff and diff3 tools, the likelihood of finding such a tool whose
calling conventions exactly match what Subversion expects is still quite
low. The key to using external differencing and merge tools with
Subversion is to use wrapper scripts, which convert the input from
Subversion into something that your specific differencing tool can
understand, and then convert the output of your tool back into a format
that Subversion expects. The following sections cover the specifics of
those expectations.

\subsection{External diff}\label{svn.advanced.externaldifftools.diff}

Subversion calls external diff programs with parameters suitable for the
GNU diff utility, and expects only that the external program will return
with a successful error code per the GNU diff definition thereof. For
most alternative diff programs, only the sixth and seventh
arguments---the paths of the files that represent the left and right
sides of the diff, respectively---are of interest. Note that Subversion
runs the diff program once per modified file covered by the Subversion
operation, so if your program runs in an asynchronous fashion (or is
``backgrounded''), you might have several instances of it all running
simultaneously. Finally, Subversion expects that your program return an
error code of 1 if your program detected differences, or 0 if it did
not---any other error code is considered a fatal error.\footnote{The GNU
  diff manual page puts it this way: ``An exit status of 0 means no
  differences were found, 1 means some differences were found, and 2
  means trouble.''}

\hyperref[svn.advanced.externaldifftools.diff.ex-1]{diffwrap.py} and
\hyperref[svn.advanced.externaldifftools.diff.ex-2]{diffwrap.bat} are
templates for external diff tool wrappers in the Python and Windows
batch scripting languages, respectively.

\protect\phantomsection\label{svn.advanced.externaldifftools.diff.ex-1}
diffwrap.py

\begin{verbatim}
#!/usr/bin/env python
import sys
import os

# Configure your favorite diff program here.
DIFF = "/usr/local/bin/my-diff-tool"

# Subversion provides the paths we need as the last two parameters.
LEFT  = sys.argv[-2]
RIGHT = sys.argv[-1]

# Call the diff command (change the following line to make sense for
# your diff program).
cmd = [DIFF, '--left', LEFT, '--right', RIGHT]
os.execv(cmd[0], cmd)

# Return an errorcode of 0 if no differences were detected, 1 if some were.
# Any other errorcode will be treated as fatal.
\end{verbatim}

\protect\phantomsection\label{svn.advanced.externaldifftools.diff.ex-2}
diffwrap.bat

\begin{verbatim}
@ECHO OFF

REM Configure your favorite diff program here.
SET DIFF="C:\Program Files\Funky Stuff\My Diff Tool.exe"

REM Subversion provides the paths we need as the last two parameters.
REM These are parameters 6 and 7 (unless you use svn diff -x, in
REM which case, all bets are off).
SET LEFT=%6
SET RIGHT=%7

REM Call the diff command (change the following line to make sense for
REM your diff program).
%DIFF% --left %LEFT% --right %RIGHT%

REM Return an errorcode of 0 if no differences were detected, 1 if some were.
REM Any other errorcode will be treated as fatal.
\end{verbatim}

\subsection{External diff3}\label{svn.advanced.externaldifftools.diff3}

Subversion invokes three-way differencing programs to perform
non-interactive merges. When configured to use an external three-way
differencing program, it executes that program with parameters suitable
for the GNU diff3 utility, expecting that the external program will
return with a successful error code and that the full file contents that
result from the completed merge operation are printed on the standard
output stream (so that Subversion can redirect them into the appropriate
version-controlled file). For most alternative merge programs, only the
ninth, tenth, and eleventh arguments, the paths of the files which
represent the ``mine'', ``older'', and ``yours'' inputs, respectively,
are of interest. Note that because Subversion depends on the output of
your merge program, your wrapper script must not exit before that output
has been delivered to Subversion. When it finally does exit, it should
return an error code of 0 if the merge was successful, or 1 if
unresolved conflicts remain in the output---any other error code is
considered a fatal error.

\hyperref[svn.advanced.externaldifftools.diff3.ex-1]{diff3wrap.py} and
\hyperref[svn.advanced.externaldifftools.diff3.ex-2]{diff3wrap.bat} are
templates for external three-way differencing tool wrappers in the
Python and Windows batch scripting languages, respectively.

\protect\phantomsection\label{svn.advanced.externaldifftools.diff3.ex-1}
diff3wrap.py

\begin{verbatim}
#!/usr/bin/env python
import sys
import os

# Configure your favorite three-way diff program here.
DIFF3 = "/usr/local/bin/my-diff3-tool"

# Subversion provides the paths we need as the last three parameters.
MINE  = sys.argv[-3]
OLDER = sys.argv[-2]
YOURS = sys.argv[-1]

# Call the three-way diff command (change the following line to make
# sense for your three-way diff program).
cmd = [DIFF3, '--older', OLDER, '--mine', MINE, '--yours', YOURS]
os.execv(cmd[0], cmd)

# After performing the merge, this script needs to print the contents
# of the merged file to stdout.  Do that in whatever way you see fit.
# Return an errorcode of 0 on successful merge, 1 if unresolved conflicts
# remain in the result.  Any other errorcode will be treated as fatal.
\end{verbatim}

\protect\phantomsection\label{svn.advanced.externaldifftools.diff3.ex-2}
diff3wrap.bat

\begin{verbatim}
@ECHO OFF

REM Configure your favorite three-way diff program here.
SET DIFF3="C:\Program Files\Funky Stuff\My Diff3 Tool.exe"

REM Subversion provides the paths we need as the last three parameters.
REM These are parameters 9, 10, and 11.  But we have access to only
REM nine parameters at a time, so we shift our nine-parameter window
REM twice to let us get to what we need.
SHIFT
SHIFT
SET MINE=%7
SET OLDER=%8
SET YOURS=%9

REM Call the three-way diff command (change the following line to make
REM sense for your three-way diff program).
%DIFF3% --older %OLDER% --mine %MINE% --yours %YOURS%

REM After performing the merge, this script needs to print the contents
REM of the merged file to stdout.  Do that in whatever way you see fit.
REM Return an errorcode of 0 on successful merge, 1 if unresolved conflicts
REM remain in the result.  Any other errorcode will be treated as fatal.
\end{verbatim}

\subsection{External merge}\label{svn.advanced.externaldifftools.merge}

Subversion optionally invokes an external merge tool as part of its
support for interactive conflict resolution. It provides as arguments to
the merge tool the following: the path of the unmodified base file, the
path of the ``theirs'' file (which contains upstream changes), the path
of the ``mine'' file (which contains local modifications), the path of
the file into which the final resolved contents should be stored by the
merge tool, and the working copy path of the conflicted file (relative
to the original target of the merge operation). The merge tool is
expected to return an error code of 0 to indicate success, or 1 to
indicate failure.

\hyperref[svn.advanced.externaldifftools.merge.ex-1]{mergewrap.py} and
\hyperref[svn.advanced.externaldifftools.merge.ex-2]{mergewrap.bat} are
templates for external merge tool wrappers in the Python and Windows
batch scripting languages, respectively.

\protect\phantomsection\label{svn.advanced.externaldifftools.merge.ex-1}
mergewrap.py

\begin{verbatim}
#!/usr/bin/env python
import sys
import os

# Configure your favorite merge program here.
MERGE = "/usr/local/bin/my-merge-tool"

# Get the paths provided by Subversion.
BASE   = sys.argv[1]
THEIRS = sys.argv[2]
MINE   = sys.argv[3]
MERGED = sys.argv[4]
WCPATH = sys.argv[5]

# Call the merge command (change the following line to make sense for
# your merge program).
cmd = [MERGE, '--base', BASE, '--mine', MINE, '--theirs', THEIRS,
              '--outfile', MERGED]
os.execv(cmd[0], cmd)

# Return an errorcode of 0 if the conflict was resolved; 1 otherwise.
# Any other errorcode will be treated as fatal.
\end{verbatim}

\protect\phantomsection\label{svn.advanced.externaldifftools.merge.ex-2}
mergewrap.bat

\begin{verbatim}
@ECHO OFF

REM Configure your favorite merge program here.
SET MERGE="C:\Program Files\Funky Stuff\My Merge Tool.exe"

REM Get the paths provided by Subversion.
SET BASE=%1
SET THEIRS=%2
SET MINE=%3
SET MERGED=%4
SET WCPATH=%5

REM Call the merge command (change the following line to make sense for
REM your merge program).
%MERGE% --base %BASE% --mine %MINE% --theirs %THEIRS% --outfile %MERGED%

REM Return an errorcode of 0 if the conflict was resolved; 1 otherwise.
REM Any other errorcode will be treated as fatal.
\end{verbatim}

\section{Summary}\label{svn.customization.summary}

Sometimes there\textquotesingle s a single right way to do things;
sometimes there are many. Subversion\textquotesingle s developers
understand that while the majority of its exact behaviors are acceptable
to most of its users, there are some corners of its functionality where
such a universally pleasing approach doesn\textquotesingle t exist. In
those places, Subversion offers users the opportunity to tell it how
\emph{they} want it to behave.

In this chapter, we explored Subversion\textquotesingle s runtime
configuration system and other mechanisms by which users can control
those configurable behaviors. If you are a developer, though, the next
chapter will take you one step further. It describes how you can further
customize your Subversion experience by writing your own software
against Subversion\textquotesingle s libraries.

